La consecuencia conceptual se concentra en la articulación de \(K\),
\(\alpha\) y la sección electrónica. La reconstrucción reversible de
\(K\) conserva datos que intervienen tanto en la normalización posicional
como en la incidencia de frontera. La composición electrónica utiliza las
coordenadas previamente construidas y mantiene su centro, sus transportes
y su realización dimensional. La prolongación excepcional conserva otras
relaciones del mismo estado bajo representaciones sucesivas. Las conexiones
entre estos objetos forman, por tanto, parte del resultado junto con los
valores de sus coordenadas.

Esta organización determina también el orden de la contrastación. Cuando
una derivación fija un valor sin recibirlo como argumento, la medición
somete a prueba la consecuencia obtenida. La correspondencia con un
observable físico exige, además, identificar la magnitud, la normalización
y la realización dimensional empleadas. En ese sentido preciso, la tesis
examina el tránsito de una descripción parametrizada a la explicación de
las relaciones que determinan sus parámetros. El alcance de cada resultado
es el de las demostraciones expuestas: estructura, ecuación y valor
constituyen aspectos inseparables de la determinación que se contrasta.

La unicidad del registro conserva dos censos prospectivos y sus
operaciones propias (apartado~\ref{act21:sec:relacion-censos}).
La cadena \(40320\to21902\to19446\to79\to1\) culmina en el
superviviente heterotípico de la cara excepcional; la cadena
\(196\cdot197\cdot196\to331\to2\to1\) culmina en la selección
orientada desde catálogos regionales. Las poblaciones intermedias
conservan el significado de cada filtro. La composición de las salidas
seleccionadas con la inversión integral articula determinación
prospectiva y recuperación reversible.
